\section{Held-out confirmation}
\label{sec:heldout}

We applied the pre-fixed \ctrl{shrinking\_Q4\_narrow} to
4 specs $\times$ seeds 100--109 $=$ \textbf{40 instances}: 960 runs, zero
failures, zero floor caps. \emph{The configuration was not reselected.}

\begin{table}[!ht]
  \centering
  % 이 파일은 scripts/make_tables.py 가 생성한다. **손으로 수정하지 않는다.**
% 표: held-out absolute
% raw: headroom_challenge-heldout_step_size_fixed_b8_9a18b6e9.jsonl
% raw 의 SHA-256 은 paper/evidence_map.md 에 있다.
% 공개 저장소에는 raw 가 없다. results/public/*.csv 에서 --public-dir 로 만들면
% 이 파일과 바이트 단위로 같아진다 (scripts/verify_public_results.py).
\begin{tabular}{lrrrr}
\toprule
Controller & $\kappa = 10^{3}$ & $\kappa = 10^{4}$ & $\kappa = 10^{5}$ & $\kappa = 10^{6}$ \\
\midrule
\texttt{best\_static} & 12.741 & 8.825 & 8.980 & 8.724 \\
\texttt{best\_open\_loop} & 13.656 & 9.332 & 9.251 & 9.086 \\
\texttt{heuristic} & 12.740 & 8.825 & 8.980 & 8.755 \\
\texttt{onestep\_narrow} & 17.032 & 10.357 & 9.879 & 9.833 \\
\texttt{onestep\_absolute} & 17.036 & 10.571 & 9.882 & 9.723 \\
\texttt{committed\_Q4\_narrow} & 19.417 & 11.345 & 10.205 & 10.123 \\
\texttt{shrinking\_Q4\_narrow} & 19.911 & 11.222 & 10.186 & 10.032 \\
\bottomrule
\end{tabular}

  \caption{Absolute median $\JE$ in nats on held-out instances ($n = 10$ per
  spec). Higher is better.}
  \label{tab:heldout-absolute}
\end{table}

\begin{table}[!ht]
  \centering
  % 이 파일은 scripts/make_tables.py 가 생성한다. **손으로 수정하지 않는다.**
% 표: held-out paired delta
% raw: headroom_challenge-heldout_step_size_fixed_b8_9a18b6e9.jsonl
% raw 의 SHA-256 은 paper/evidence_map.md 에 있다.
% 공개 저장소에는 raw 가 없다. results/public/*.csv 에서 --public-dir 로 만들면
% **고정된 환경에서** 이 파일과 바이트 단위로 같아진다.
% 판본이 다른 환경에서의 비트 수준 일치는 주장하지 않는다.
\begin{tabular}{llrrrr}
\toprule
Gate & Comparison & Median & 95\% CI & $p$ & Positive \\
\midrule
A2 & shrinking $-$ best\_static & $+1.690$ & $[+1.462, +2.368]$ & ${<}0.0001$ & 40/40 \\
C2 & shrinking $-$ onestep & $+0.456$ & $[+0.254, +0.720]$ & ${<}0.0001$ & 35/40 \\
C3 & shrinking $-$ committed & $+0.010$ & $[-0.033, +0.053]$ & $0.9725$ & 21/40 \\
B & onestep\_absolute $-$ onestep\_narrow & $+0.005$ & $[-0.002, +0.026]$ & $0.5900$ & 23/40 \\
B & onestep\_wide $-$ onestep\_narrow & $-0.001$ & $[-0.041, +0.005]$ & $0.3611$ & 18/40 \\
\midrule
-- & open\_loop $-$ best\_static & $+0.395$ & $[+0.350, +0.476]$ & ${<}0.0001$ & 40/40 \\
-- & heuristic $-$ best\_static & $-0.000$ & $[-0.000, +0.000]$ & $0.7750$ & 20/40 \\
\bottomrule
\end{tabular}

  \caption{Paired differences on held-out instances ($n = 40$). Positive values
  favour the treatment. Gate verdicts:
  \gate{A1}~GO, \gate{A2}~GO, \gate{B}~redesign, \gate{C1}~undecidable,
  \gate{C2}~GO, \gate{C3}~redesign, \gate{D}~GO.}
  \label{tab:heldout-paired}
\end{table}

\begin{figure}[!ht]
  \centering
  \includegraphics[width=\textwidth]{figure1_ladder.png}
  \caption{Absolute performance across the ladder. This figure describes levels
  only. \textbf{Do not subtract bars in panel (b);} the median of paired
  differences is not linear, so the causal decomposition must be read from
  Figure~\ref{fig:planning-vs-feedback} instead.}
  \label{fig:ladder}
\end{figure}

\subsection{Multi-step lookahead beats one-step}
\label{sec:c2}

$\gate{C2} = +0.456$~nat with a CI lower bound of $+0.254$, above the GO threshold
of $0.3$. Both required conditions were met:

\begin{center}
\begin{tabular}{lr}
\toprule
Condition & Value \\
\midrule
Improvement & $+0.456$~nat, $35/40$ positive \\
Adoption rate of depth $> 1$ & 0.84 \\
Plan-depth cap reached & 0.00 \\
\bottomrule
\end{tabular}
\end{center}

No step failed to consume its quota because of a computational cap, so the
quota-ladder comparison is not compromised.

On development seeds this was conditional ($+0.251$, $p = 0.0771$). At $n = 40$ it
was promoted to GO.

\subsection{Replanning during execution}
\label{sec:c3}

$\gate{C3} = +0.010$~nat, 95\% CI $[-0.033, +0.053]$, with 21 wins, 1 tie, and 18
losses. The CI width of $0.086$~nat is an order of magnitude narrower than the
$0.906$~nat of \gate{A2}.

\begin{quote}
In the held-out results the feedback effect was $+0.010$~nat with a 95\%
confidence interval of $[-0.033, +0.053]$ that narrowly included zero. We
therefore \emph{did not observe a practically large feedback benefit}.
\end{quote}

We do not claim equivalence. Because we did not preregister an equivalence margin,
we cannot report ``the effect is zero'' or ``the effect is smaller than
$0.053$~nat'' as a test result. The two one-sided tests procedure requires bounds
fixed in advance \citep{schuirmann1987tost}, and specifying those bounds ahead of
time is the recommended practice \citep{lakens2017equivalence}.

\subsection{The action space is not the bottleneck}
\label{sec:actionspace}

\ctrl{absolute} (132 actions, $\log_{10}$ range 15.27) gains $+0.005$~nat over
\ctrl{narrow} (12 actions, range 0.95), and $\ctrl{wide} - \ctrl{narrow}$ is
$-0.001$~nat; see Table~\ref{tab:heldout-paired} for intervals. Letting damping be
chosen freely does not raise one-step performance. \emph{This is a negative result,
but it justifies the inexpensive narrow action space.}

The GO threshold for \gate{B} was $0.5$~nat and the observation is two orders of
magnitude smaller, so the verdict is \emph{redesign}. The interval includes zero,
but we make no equivalence claim, for the reason given in Section~\ref{sec:c3}.

\Needspace{6\baselineskip}
\subsection{Search cost}
\label{sec:cost}

\begin{table}[!ht]
  \centering
  % 이 파일은 scripts/make_tables.py 가 생성한다. **손으로 수정하지 않는다.**
% 표: decision-search cost
% raw: headroom_challenge-heldout_step_size_fixed_b8_9a18b6e9.jsonl
% raw 의 SHA-256 은 paper/evidence_map.md 에 있다.
% 공개 저장소에는 raw 가 없다. results/public/*.csv 에서 --public-dir 로 만들면
% 이 파일과 바이트 단위로 같아진다 (scripts/verify_public_results.py).
\begin{tabular}{lrr}
\toprule
Controller & Decision-search GE & Relative to budget \\
\midrule
\texttt{best\_static} & 0 & -- \\
\texttt{best\_open\_loop} & 0 & -- \\
\texttt{heuristic} & 0 & -- \\
\texttt{onestep\_narrow} & 1,186 & $7.9\times$ \\
\texttt{onestep\_absolute} & 10,084 & $67.2\times$ \\
\texttt{committed\_Q4\_narrow} & 69,401 & $462.7\times$ \\
\texttt{shrinking\_Q4\_narrow} & 194,095 & $1,294.0\times$ \\
\bottomrule
\end{tabular}

  \caption{Median decision-search cost on held-out instances, relative to the
  $150$~GE deployment budget. The reported headroom was obtained by an oracle that
  consumed $1{,}294$ times the deployment budget; Section~\ref{sec:cheap} discusses
  what that cost gap implies.}
  \label{tab:heldout-cost}
\end{table}
